\chapter{FUTURE WORK}

\section{Introduction}

While the current system achieves strong results for binary waste classification, several natural extensions would enhance its capability, coverage, and real-world impact. This chapter outlines the most promising directions for future development.

\section{Multi-Category Waste Classification}

The current system classifies waste into two categories (Bio / Non-Bio). Real-world waste management requires finer-grained segregation:
\begin{itemize}
    \item \textbf{Paper and Cardboard} (recyclable, separate processing from plastic)
    \item \textbf{Glass} (requires dedicated recycling stream)
    \item \textbf{Metal} (high-value recyclable)
    \item \textbf{Plastic} (multiple sub-types: PET, HDPE, PVC)
    \item \textbf{Electronic Waste} (hazardous, specialised disposal)
    \item \textbf{Organic/Food Waste} (composting stream)
\end{itemize}

Extending to 6--10 categories would require: (i) a larger annotated dataset (at least 10,000 images per category), (ii) class-weighted training to handle natural class imbalance, and (iii) re-evaluation of the BiFPN configuration for the expanded label space.

\section{Hardware Integration}

The current system is software-only. Integration with physical bin actuation hardware represents the next step toward a complete automated segregation system:
\begin{itemize}
    \item \textbf{Servo-motor-actuated bin lids:} Detection output triggers the appropriate bin compartment to open.
    \item \textbf{Conveyor belt control:} Classification signal redirects waste items on a sorting conveyor.
    \item \textbf{Raspberry Pi / Jetson deployment:} Model compression (quantisation, pruning) to run inference on embedded hardware.
    \item \textbf{Arduino integration:} PySerial communication from inference script to Arduino microcontroller for actuator control.
\end{itemize}

\section{Mobile Application}

A mobile application would make the system accessible for household and institutional use without requiring a laptop:
\begin{itemize}
    \item \textbf{Android/iOS app:} Camera capture $\rightarrow$ on-device inference $\rightarrow$ classification result displayed.
    \item \textbf{Model optimisation:} Export to ONNX or TensorFlow Lite for mobile runtime, applying quantisation (INT8) to reduce model size below 5\,MB.
    \item \textbf{Offline capability:} On-device inference ensures operation without internet connectivity.
\end{itemize}

\section{Continual Learning}

Static models degrade over time as new waste types emerge and visual distributions shift. A continual learning pipeline would:
\begin{itemize}
    \item Collect uncertain predictions (low-confidence detections) from deployed systems.
    \item Queue them for human review and annotation.
    \item Periodically retrain the model on the expanded dataset with new examples.
    \item Deploy updated weights without requiring full retraining from scratch.
\end{itemize}

This approach maintains model accuracy over time without catastrophic forgetting of previously learned categories.

\section{Fill Level and Route Optimisation Integration}

The classification system can be integrated with fill-level monitoring and smart city logistics:
\begin{itemize}
    \item \textbf{Fill monitoring:} Ultrasonic or weight sensors track bin capacity for each waste category separately.
    \item \textbf{Cloud dashboard:} Real-time bin status map with collection priority indicators.
    \item \textbf{Route optimisation:} Collection routes dynamically adjusted based on which bins require emptying, reducing fuel consumption and collection frequency.
\end{itemize}

Research has shown that such systems can reduce collection route distances by up to 32\% \cite{iot_ai2024}, representing significant operational savings when fed by accurate classification.

\section{Model Distillation for Edge Deployment}

The current 9.4M parameter model, while compact, may be too large for microcontroller-class devices. Knowledge distillation can produce a smaller student model:
\begin{itemize}
    \item \textbf{Teacher model:} Current YOLO11s + BiFPN + WIoU (94.39\% mAP)
    \item \textbf{Student model:} YOLO11n-scale with distilled feature representations
    \item \textbf{Target:} $<$3M parameters, $<$5\,MB, maintaining $>$90\% mAP@0.5
\end{itemize}

\section{Dataset Expansion and Diversity}

The current 3,000-image dataset, while sufficient for binary classification, could be expanded:
\begin{itemize}
    \item Increase to 10,000+ raw images for improved generalisation.
    \item Include nighttime and low-light images for 24-hour deployment scenarios.
    \item Add images from diverse geographic regions and cultural waste contexts.
    \item Include wet, soiled, and degraded waste items that are challenging edge cases.
\end{itemize}

\section{Explainability and Transparency}

For public-facing deployment, model explainability would increase user trust:
\begin{itemize}
    \item \textbf{Grad-CAM visualisation:} Heat maps showing which image regions influenced the classification decision.
    \item \textbf{Confidence calibration:} Ensuring that reported confidence scores accurately reflect true classification probability.
    \item \textbf{Uncertainty estimation:} Flagging low-confidence predictions for human review rather than automated action.
\end{itemize}
